\section{Edge Resource Management}
\label{sec:edge_sched}

While the RAN resource manager prioritizes LC requests during uplink transmission, the edge server must allocate compute resources to meet SLO deadlines without any coordination with the RAN.
The key challenge is estimating each request's remaining time budget by determining the network latency already consumed (uplink), the future network latency (downlink), and the expected processing time.
\modified{This section presents our approach that leverages timing signals naturally generated by applications during their request-response lifecycle, which can be captured via a lightweight API (\autoref{tbl:smec-api}).}

\begin{table}[]
\setlength{\tabcolsep}{10pt}
\renewcommand{\arraystretch}{1.0}
\centering
\footnotesize
\begin{tabular}{@{}ll@{}}
\toprule
\textbf{API Call} & \textbf{Purpose} \\ \midrule
\texttt{request\_sent(req\_data*)} & Report new request sent \\
\texttt{request\_arrived(req\_data*)} & Report new request arrival \\
\texttt{processing\_started(req\_id)} & Report processing start \\
\texttt{processing\_ended(req\_id)} & Report processing completion \\
\texttt{response\_sent(resp\_data*)} & Report response transmission \\ 
\texttt{response\_arrived(resp\_data*)} & Report response arrival \\\bottomrule
\end{tabular}
\vspace{-0.5em}
\caption{\sys{} API for network and processing time estimation.}
\label{tbl:smec-api}
\end{table}


\subsection{Estimating Network Latency}
\label{sec:reference-time}

Consider a request with a \ms{100} SLO that has spent \ms{60} in uplink transmission. 
The edge server has only \ms{40} remaining to complete processing and downlink transmission.
Without accurate timing information about both consumed uplink delays and future downlink delays, the edge server cannot distinguish between urgent requests approaching their deadlines and requests with ample time remaining.

However, measuring per-request network latency is challenging: the edge server only observes when requests arrive locally but lacks visibility into when they were originally sent from the client, the uplink delays they experienced, or the downlink delays that responses will experience.

\mypara{Possible approach.}
A straightforward approach would piggyback a sending timestamp in each request, enabling the server to compute transmission delay upon reception.
Unfortunately, this requires precise time synchronization between the UE and server, which is infeasible in MEC environments.
Network Time Protocol (NTP)~\cite{ntp-rfc5905} introduces synchronization errors ranging from tens to hundreds of milliseconds, which are larger than the time budget of LC applications.
PTP~\cite{ptp-std} is also unsuitable because it assumes network delays are symmetric.
5G networks exhibit inherent asymmetry where uplink latency is both higher and more variable than downlink latency due to protocol design and UE energy constraints. 

\begin{figure}[]
    \centering
    \includegraphics[width=0.4\textwidth]{figures/time_reference.pdf}
    \vspace{-0.5em}
    \caption{Probing-based network latency estimation. 
            Red arrows indicate probing protocol packets while blue arrows indicate application requests and responses.}
    \label{fig:reference-time}
    \vspace{-1em}
\end{figure}

\mypara{Our approach.}
We leverage our observation that while uplink latency in 5G networks is highly variable, downlink latency remains consistently stable (\autoref{fig:motive-decompose}).
We exploit the stability of downlink through a lightweight \emph{probing-based network latency estimation} protocol that establishes timing references using the stable downlink path.

As illustrated in \autoref{fig:reference-time}, the client device periodically sends probe packets to the server, which responds with ACKs over the stable downlink.
To implement this protocol, \sys{} runs a per-UE \emph{probing daemon} on the client side (shown in~\autoref{fig:architecture}) and integrates the corresponding functionality into the edge resource manager on the server side.
When an application sends a request, it reports the request to the timing daemon via the \texttt{request\_sent} API (\autoref{tbl:smec-api}), which measures $t_{ack-req}$ (the time elapsed since receiving the latest ACK) and inserts this timing metadata into the request payload.
Upon receiving the request, the server computes $T_{ack-req}$ (the time difference between sending the most recent ACK and receiving the current request).
The stable downlink timing creates a \emph{parallelogram} relationship (grey region in \autoref{fig:reference-time}), allowing the server to estimate network latency as $T_{ack-req} - t_{ack-req}$.

However, if the response size is significantly larger than the ACK size, there will be a gap in downlink transmission latency between the two (\autoref{fig:motive-decompose}). 
To compensate for this difference, when the server sends a response, it computes the elapsed time since it sent the last ACK ($T_{ack-resp}$) and reports this value to the client daemon as part of the response. 
Upon receiving the response, the client daemon uses this information to compute the time elapsed since it received the last ACK ($t_{ack-resp}$) and calculates a compensation factor ($t_{comp}$), which it reports to the server as part of the next probe. 
The client daemon maintains this compensation factor separately for each application. 
Using this correction factor, the server estimates the network latency while accounting for response size differences:
\begin{equation}
    t_{network} = T_{ack-req} - t_{ack-req} + t_{comp}
\end{equation}

To handle packet losses, each probe-ACK exchange carries a unique ID that both endpoints use to synchronize on the most recent successful exchange.
\modified{The design incurs minimal overhead: the client sends small probe/ACK packets (\textless{}\byte{100}) every few seconds, but \emph{only while the UE is actively serving LC traffic}.
When the UE is idle, the probing daemon pauses, avoiding interference with the UE's power-saving mechanism (\eg Discontinuous Reception (DRX)).}

This approach extends to distributed scenarios where request initiators and response receivers are different devices (\eg smart stadium with camera initiators and audience UE receivers).
In such cases, the initiator sends probe packets for network latency estimation while receivers sends probe packets to report compensation factors, enabling accurate end-to-end latency estimation across the distributed path.

\subsection{Estimating Remaining Time Budget}
\label{sec:deadline-estimation}

Given the network latency estimates, the resource manager now needs to predict each request's \emph{processing time} to compute an accurate remaining time budget.
However, accurately predicting processing time while maintaining practical deployability presents challenges.
Applications exhibit variable processing times due to workload change and resource contention, making accurate prediction difficult.

\mypara{Possible approaches.}
One approach would be to infer processing time from system-wide signals like CPU utilization, scheduling frequency, and hardware performance counters (\eg cache misses, memory footprint), but such methods suffer from either high inference overhead or poor accuracy, as shown by existing work~\cite{liu2023intelligent, peng2024ufo}.
Alternatively, one could extensively instrument applications to measure real-time workload information (\eg Caladan~\cite{caladan:osdi20}); while this obtains precise timing information, it requires significant application modifications, limiting practicality.

\mypara{Our approach.}
Building on the client-side request tracking for network latency estimation (\autoref{sec:reference-time}), we extend this to the server side where applications report key lifecycle events (\eg request arrival, processing start/end) to the resource manager via the \sys{} API (\autoref{tbl:smec-api}).
By tracking these events, the resource manager estimates processing delays without requiring detailed application knowledge or extensive instrumentation.

The resource manager tracks two key metrics: waiting time ($t_{wait}$) defined as the time elapsed from request arrival until processing begins, and processing time using the median of last $R$ requests as a robust predictor ($t_{process}$).
While this median-based approach is simple and may introduce some prediction error, it performs well in practice (\autoref{sec:eval-micro-benchmarks-network-and-processing-time-estimation}) while minimizing application modifications.
The manager can then compute the remaining time budget as:
\begin{equation}
    t^{edge}_{budget} = SLO - (t_{network} + t_{wait} + t_{process})
\end{equation}


\subsection{Deadline-aware Proactive Edge Resource Scheduling}
\label{sec:proactive-resource-allocation}

\begin{figure}[]
    \centering
    \begin{subfigure}[b]{0.48\columnwidth}
        \centering
        \includegraphics[width=\textwidth]{figures/resource-latency-relationship/cpu/latency_cpu.pdf}
        \caption{CPU-based task}
        \label{fig:latency-cpu}
    \end{subfigure}
    \hfill
    \begin{subfigure}[b]{0.48\columnwidth}
        \centering
        \includegraphics[width=\textwidth]{figures/resource-latency-relationship/gpu/latency_priority.pdf}
        \caption{GPU-based task}
        \label{fig:latency-gpu}
    \end{subfigure}
    \vspace{1em}
    \tightcaption{Relationship between compute resources allocation and process latency.}
    \label{fig:latency-resource}
\end{figure}

\modified{The resource manager determines urgency based on each request's remaining time budget.
We define an urgency threshold as a fraction $\tau$ of the application's SLO (default $\tau=0.1$), marking a request as urgent when $t^{edge}_{budget} < \tau \times SLO$.
When a request becomes urgent, the manager \emph{proactively} allocates additional resources to prevent SLO violations.

\sys{} focuses on CPU and GPU management, which constitute the primary compute bottlenecks in edge servers.
CPU cores are heavily contended among multiple applications, while GPUs accelerate ML inference and video processing.
\autoref{alg:edge-proactive-sched} in \autoref{appnd:edge-scheduling-algo} summarizes our deadline-aware scheduling approach.}


\mypara{CPU management.}
Our CPU manager partitions cores across applications using CPU affinity to ensure isolation and predictable performance.
Based on the observation that increasing CPU core allocation reduces processing latency (\autoref{fig:latency-cpu}), the manager allocates additional cores to urgent requests when they risk missing deadlines.
This policy is most effective when an application can parallelize request processing (\eg via multi-threaded execution).
To avoid thrashing from frequent core reallocations, the manager enforces a brief cool-down period (\eg \ms{100}) after each allocation, and assigns another core only if requests still risk missing deadlines after the cool-down.
This mechanism prevents wasteful oscillations while ensuring urgent requests receive the resources they need.

For \emph{reclamation}, we use average CPU utilization rather than urgency signals.
Urgency-based reclamation creates instability: removing a single core from a latency-critical application can cause an abrupt shift from meeting deadlines to missing many, leading to scheduler thrashing.
Instead, when an application's utilization falls below a threshold (\eg 60\%), the manager safely reclaims cores, providing stable resource management without oscillatory behavior.

\mypara{GPU management.}
Commercial MEC offerings~\cite{www-outpost,www-azure-stack-edge} typically deploy inference-optimized GPUs such as NVIDIA L4~\cite{nvidia-l4} and T4~\cite{nvidia-t4}.
These devices lack hardware-level partitioning capabilities like Multi-Instance GPU (MIG)~\cite{nvidia-mig}.
\modified{Even when hardware partitioning is available, it only supports static allocation at application launch and cannot be dynamically adjusted during runtime.

To enable dynamic GPU scheduling, we leverage CUDA stream priorities, inspired by Orion~\cite{strati2024orion}.
NVIDIA's Multi-Process Service (MPS)~\cite{nvidia-mps} enables multiple applications to share the GPU while maintaining a unified priority hierarchy across their CUDA streams.
When kernels from multiple applications contend for GPU resources, those launched from higher-priority streams receive preferential scheduling.
This mechanism allows the resource manager to influence GPU scheduling decisions without requiring hardware partitioning or GPU driver modifications.

We exploit this capability by assigning stream priorities at request granularity based on urgency.
As shown in \autoref{fig:latency-gpu}, higher stream priorities reduce processing latency under contention.
The resource manager assigns higher-priority streams to requests whose expected processing time closely matches their remaining time budget, while requests with slack use lower-priority tiers.
This ensures urgent requests receive preferential GPU access without starving other workloads.

\mypara{Early drop: handling overly urgent requests.}
When a request becomes overly urgent (\eg $t^{edge}_{budget}\leq 0$) due to excessive network or queuing delays, no amount of compute resources can recover the already-elapsed time.
Processing such requests wastes resources that could serve other requests still capable of meeting their deadlines.
Therefore, when the edge server operates under load, the resource manager immediately drops overly urgent requests (``early drop''), redirecting resources to requests with viable time budgets.}